\justifying
\NSFCBodyText

\subsubsection{年度研究计划}

本课题按一年研究周期组织，以下月份均相对于正式立项时间计算；起止年月以\textbf{【待补 F-EDU-03：正式指南或任务书确定的执行期】}为准。按照风险定义先行、评价校准承接、干预比较收束的依赖关系推进，避免在标注口径尚未稳定时同时扩展多个场景。

\textbf{【暂定 第1—3个月】}，完成两类教育任务的边界设计、公开资料与使用条件核对，建立隔离工具环境和轨迹采集方式。围绕风险类别、首次可观察位置和教学支持要求编制标注手册，以小样本预实验检查分歧、运行成本和有效风险比例，据此调整任务家族数量与主要比较。阶段交付为任务规范、标注样例、预实验分析和可复现的环境配置。

\textbf{【暂定 第4—6个月】}，完成基础任务与变体构建、独立标注及数据划分，建立结果评价、通用轨迹评价和教育情境评价的对照。开展教育约束消融，使用校准集确定监测工作点与效用要求，并在读取测试结果之前明确主要终点和统计方案。同步撰写方法部分，阶段交付为评测数据初版、证据规则、评价脚本和冻结的分析计划。

\textbf{【暂定 第7—9个月】}，开展留出任务的轨迹评测与分支干预比较，报告风险、误报、正常任务完成和人工负担。针对主要失败类别补充机制分析；在核心任务完成且资源允许时开展双智能体传播压力测试。本阶段形成主要结果论文初稿，阶段交付为主要实验报告、监测原型及失败案例说明，不以追求单一最好分数替换预定比较。

\textbf{【暂定 第10—12个月】}，完成授权条件下的高校案例分析；若真实数据或场景未获许可，则采用明确标注的近真实沙箱案例。整理风险表征、评价器与监测方法的适用范围，完成论文和典型案例报告，清理可公开数据与代码，形成可复核的成果包。论文发表或录用进度以期刊实际审理为准，结题认定口径需依据正式任务书确认。

\subsubsection{预期研究结果}

依据项目输入材料，目标为发表SCI、SSCI或CSSCI论文不少于1篇，并完成1份典型案例研究报告。围绕这两项成果，拟形成教育智能体风险分类与标注规范1套、可复用的任务与轨迹评测集1套，以及监测原型和配套评价脚本。论文、基准和原型共享同一实验主线，避免将同一工作拆成彼此不相关的多项承诺。

评测集将同时提供正常任务、受控风险任务、人工判断依据、数据划分与运行设置，公开范围服从原材料许可和授权约定。研究报告不仅呈现方法收益，也记录无增益、误告警和过度干预等反例，给出需要人工复核的条件。软件著作权和额外高水平论文作为资源允许时的争取成果，不作为扩大主研究规模的理由。
